\abstract{
  We construct the Lorentzian completion of the spectral entropy framework
  in which the renormalized variation of a Laplace-type operator generates
  an infrared Einstein response.
  While the original formulation was elliptic and sufficient to derive
  the effective field equations and their thermodynamic interpretation,
  it did not address causal propagation.
  We define a Lorentzian spectral action via the real part of the
  hyperbolic determinant, equivalently through a Schwinger--Keldysh
  prescription, ensuring a real and retarded quadratic kernel.

  Linearization about Minkowski spacetime shows that, in the regime
  $R\ell_\chi^2 \ll 1$ and $\omega \ll \ell_\chi^{-1}$,
  the theory propagates exactly two transverse-traceless tensor modes of
  helicity $\pm2$.
  Scalar and vector sectors remain non-dynamical in the infrared, and
  the additional pole generated by higher-derivative corrections lies at
  $k^2 \sim \ell_\chi^{-2}$, beyond current observational reach.
  The polarization content therefore coincides with that of general relativity
  in the accessible domain.
  When the vector sector is sourced by a projected momentum density it obeys an elliptic gravitomagnetic
  constraint, fixing the frame-dragging potential $g_{0i}$ without an additional propagating degree of freedom,
  conditionally on the projected matter current.

  Beyond the leading Einstein sector, the ratio of Seeley--DeWitt
  coefficients fixes a universal $k^4$ correction to the dispersion relation,
  \[
    \omega^2 = c^2 k^2 - \gamma \ell_\chi^2 k^4 + \mathcal{O}(k^6),
  \]
  with $\gamma = 1/(180\zeta)$ and $\zeta=\mathcal{O}(1)$ encoding scheme-dependent normalization factors.
  In the natural spectral scheme one finds $\gamma=1/30$.
  This structure maps directly onto the parametrizations used in
  gravitational-wave catalogues and yields a bound on the pre-geometric
  scale $\ell_\chi$ from interferometric data.

  The Lorentzian completion thus embeds the spectral entropy dynamics
  in a fully causal framework and establishes gravitational waves as
  infrared probes of the underlying pre-geometric scale.
}
